\chapter{Vectoren en rechten in het vlak}\label{ch:g11:vect}

\omterm{def:g11:vect:vector}{Vectoren} leggen verplaatsingen vast; rechten zijn hun banen. Het belangrijke
nieuwe gereedschap van dit hoofdstuk is de \emph{\omterm{def:g11:vect:det}{determinant}}, één getal dat
beslist of twee \omterm{def:g11:vect:vector}{vectoren} \omterm{def:g11:vect:collinear}{collineair} zijn --- en dat van daaruit \omterm{def:g10:algebra:equation}{vergelijkingen}
van rechten voortbrengt en vragen over collineariteit en evenwijdigheid met
zuiver rekenwerk beslecht. Dezelfde ideeën gaan in \cref{ch:g12:space} naar de
ruimte.

\section{Vectoren: een opfrissing}

\begin{definition}[Vector, coördinaten]\label{def:g11:vect:vector}
Een \emph{vector}\index{vector} $\vec u$ stelt een verplaatsing voor: alle
pijlen met dezelfde richting, zin en lengte stellen dezelfde vector voor. In
een \omterm{def:g10:coordgeom:system}{assenstelsel} noteert $\vec u\,(x, y)$ de verplaatsing ($x$ horizontaal,
$y$ verticaal); voor punten $A(x_A, y_A)$ en $B(x_B, y_B)$ is
\[
\vect{AB}\,\bigl(x_B - x_A,\ y_B - y_A\bigr).
\]
Vectoren tellen coördinaat per coördinaat op, en $\lambda\vec u$ heeft
\omterm{def:g10:coordgeom:system}{coördinaten} $(\lambda x, \lambda y)$.
\end{definition}

\begin{proposition}[Midden]\label{prop:g11:vect:midpoint}
Het \omterm{prop:g10:coordgeom:midpoint}{midden} $M$ van het lijnstuk $[AB]$ heeft \omterm{def:g10:coordgeom:system}{coördinaten}
$\left(\frac{x_A + x_B}{2},\ \frac{y_A + y_B}{2}\right)$.
\end{proposition}

\begin{proof}
$M$ wordt gekenmerkt door $\vect{AM} = \frac12\vect{AB}$: coördinaat per
coördinaat geeft dat $x_M - x_A = \frac12(x_B - x_A)$, dus
$x_M = \frac{x_A + x_B}{2}$, en net zo voor $y_M$.
\end{proof}

\section{Collineariteit en de determinant}

\begin{definition}[Collineaire vectoren]\label{def:g11:vect:collinear}
Twee \omterm{def:g11:vect:vector}{vectoren} heten \emph{collineair}\index{collineaire vectoren} wanneer de
ene een scalair veelvoud van de andere is: $\vec v = \lambda \vec u$ voor een
zekere $\lambda \in \R$ (of $\vec u = \vec 0$). \omterm{def:g11:seq:geometric}{Meetkundig}: ze hebben
dezelfde richting, of een van beide is nul.
\end{definition}

\begin{definition}[Determinant]\label{def:g11:vect:det}
De \emph{determinant}\index{determinant} van $\vec u\,(x, y)$ en
$\vec v\,(x', y')$ is
\[
\det(\vec u, \vec v) =
\begin{vmatrix} x & x' \\ y & y' \end{vmatrix}
= x y' - x' y .
\]
\end{definition}

\begin{theorem}[Criterium voor collineariteit]\label{thm:g11:vect:collinearity}
Twee \omterm{def:g11:vect:vector}{vectoren} $\vec u$ en $\vec v$ zijn \omterm{def:g11:vect:collinear}{collineair} dan en slechts dan als
$\det(\vec u, \vec v) = 0$.
\end{theorem}

\begin{proof}
($\Rightarrow$) Is $\vec v = \lambda\vec u$, dan is $x' = \lambda x$ en
$y' = \lambda y$, dus $\det(\vec u, \vec v) = x(\lambda y) - (\lambda x)y
= 0$; is $\vec u = \vec 0$, dan is de \omterm{def:g11:vect:det}{determinant} vanzelf $0$.

($\Leftarrow$) Stel $xy' - x'y = 0$ en $\vec u \neq \vec 0$, zeg $x \neq 0$
(het geval $y \neq 0$ verloopt symmetrisch). Stel $\lambda = \frac{x'}{x}$.
Dan is $x' = \lambda x$ per constructie, en de betrekking
$xy' = x'y = \lambda x y$ geeft na deling door $x \neq 0$ dat
$y' = \lambda y$. Bijgevolg is $\vec v = \lambda \vec u$.
\end{proof}

\begin{example}\label{ex:g11:vect:collinear}
$\vec u\,(3, -2)$ en $\vec v\,(-6, 4)$:
$\det = 3 \times 4 - (-6) \times (-2) = 12 - 12 = 0$, dus \omterm{def:g11:vect:collinear}{collineair} (en
inderdaad is $\vec v = -2\vec u$). $\vec u\,(3, -2)$ en $\vec w\,(2, 1)$:
$\det = 3 + 4 = 7 \neq 0$, dus niet \omterm{def:g11:vect:collinear}{collineair}.
\end{example}

\begin{method}[Drie punten op één rechte]\label{met:g11:vect:alignment}
$A$, $B$, $C$ liggen op één rechte dan en slechts dan als $\vect{AB}$ en
$\vect{AC}$ \omterm{def:g11:vect:collinear}{collineair} zijn: bereken beide \omterm{def:g11:vect:vector}{vectoren} en ga
$\det\bigl(\vect{AB}, \vect{AC}\bigr) = 0$ na.
\end{method}

\section{Cartesische vergelijkingen van rechten}

\begin{definition}[Richtingsvector]\label{def:g11:vect:direction}
Een \emph{richtingsvector}\index{richtingsvector} van een rechte $d$ is elke
\omterm{def:g11:vect:vector}{vector} $\vec u \neq \vec 0$ waarvoor $\vect{AB}$ \omterm{def:g11:vect:collinear}{collineair} is met $\vec u$
voor alle punten $A, B$ van $d$.
\end{definition}

\begin{theorem}[Cartesische vergelijking]\label{thm:g11:vect:cartesian}
Een rechte met \omterm{def:g11:vect:direction}{richtingsvector} $\vec u\,(-b, a)$ heeft een \omterm{def:g10:algebra:equation}{vergelijking} van
de vorm
\[
ax + by + c = 0
\qquad (a, b) \neq (0, 0).
\]
Omgekeerd beschrijft elke zulke \omterm{def:g10:algebra:equation}{vergelijking} een rechte met \omterm{def:g11:vect:direction}{richtingsvector}
$\vec u\,(-b, a)$.
\end{theorem}

\begin{proof}
Zij $d$ de rechte door $A(x_A, y_A)$ met richting $\vec u\,(-b, a)$. Een
punt $M(x, y)$ ligt op $d$ precies wanneer $\vect{AM}$ \omterm{def:g11:vect:collinear}{collineair} is met
$\vec u$, dus (\cref{thm:g11:vect:collinearity}) wanneer
\[
0 = \det\bigl(\vect{AM}, \vec u\bigr)
= (x - x_A)\,a - (-b)(y - y_A)
= ax + by - (a x_A + b y_A),
\]
een \omterm{def:g10:algebra:equation}{vergelijking} van de aangekondigde vorm, met $c = -(a x_A + b y_A)$.
Omgekeerd zijn de oplossingen van $ax + by + c = 0$, met bijvoorbeeld
$b \neq 0$, de punten $\left(x, -\frac{a x + c}{b}\right)$: twee ervan van
elkaar aftrekken toont dat elke koorde \omterm{def:g11:vect:collinear}{collineair} is met $(-b, a)$, en er
bestaat altijd minstens één oplossing --- het gaat dus om de rechte door dat
punt met richting $(-b, a)$.
\end{proof}

\begin{method}[Rechte door twee punten]\label{met:g11:vect:twopoints}
Om een \omterm{def:g10:algebra:equation}{vergelijking} van de rechte $(AB)$ te vinden: bereken de
\omterm{def:g11:vect:direction}{richtingsvector} $\vect{AB}\,(x_B - x_A, y_B - y_A)$; leg hem naast $(-b, a)$,
dus neem $a = y_B - y_A$ en $b = -(x_B - x_A)$; bepaal $c$ door de
\omterm{def:g10:coordgeom:system}{coördinaten} van $A$ in te vullen.
\end{method}

\begin{example}\label{ex:g11:vect:twopoints}
De rechte door $A(1, 2)$ en $B(4, 3)$: $\vect{AB}\,(3, 1)$, dus $a = 1$,
$b = -3$, en de \omterm{def:g10:algebra:equation}{vergelijking} is $x - 3y + c = 0$. $A$ invullen geeft
$1 - 6 + c = 0$, dus $c = 5$. \omterm{def:g10:algebra:equation}{Vergelijking}: $x - 3y + 5 = 0$. Controle met
$B$: $4 - 9 + 5 = 0$.
\end{example}

\begin{remark}
Is $b \neq 0$, dan kun je de \omterm{def:g10:algebra:equation}{vergelijking} naar $y$ oplossen: $y = mx + p$
met \omterm{def:g10:reffunc:affine}{richtingscoëfficiënt} $m = -\frac ab$. De cartesische vorm
$ax + by + c = 0$ is algemener: ze dekt ook de verticale rechten ($b = 0$),
die geen \omterm{def:g10:reffunc:affine}{richtingscoëfficiënt} hebben.
\end{remark}

\begin{omfigure}
\begin{tikzpicture}
\begin{axis}[omaxis,
  width=8.6cm, height=5.8cm,
  xmin=-1.6, xmax=5.6, ymin=-0.9, ymax=4.4,
  xtick={1,4}, ytick={2,3}]
  \addplot[omDef, thick, domain=-1.4:5.4] {(x+5)/3};
  \fill (axis cs:1,2) circle (1.5pt)
    node[above left, font=\small] {$A$};
  \fill (axis cs:4,3) circle (1.5pt)
    node[above left, font=\small] {$B$};
  \draw[-Stealth, omThm, thick] (axis cs:1,2) -- (axis cs:2.5,2.5)
    node[midway, below right, font=\small] {$\vec u$};
\end{axis}
\end{tikzpicture}

{\small De rechte $x - 3y + 5 = 0$ van \cref{ex:g11:vect:twopoints}, haar
\omterm{def:g11:vect:direction}{richtingsvector} $\vec u = \vect{AB}$ (rood, hier met $\frac12$ geschaald), en
de twee punten waarmee de \omterm{def:g10:algebra:equation}{vergelijking} gebouwd werd.}
\end{omfigure}

\section{Parametervoorstelling}

\begin{definition}[Parametervoorstelling]\label{def:g11:vect:parametric}
De rechte door $A(x_A, y_A)$ met richting $\vec u\,(\alpha, \beta)$ is de
verzameling van de punten
\[
M\bigl(x_A + t\alpha,\ y_A + t\beta\bigr), \qquad t \in \R .
\]
Het getal $t$ is de \emph{parameter}: elke waarde van $t$ levert één punt van
de rechte op, alsof je haar met constante snelheid $\vec u$ doorloopt.
\end{definition}

\begin{example}
De rechte van \cref{ex:g11:vect:twopoints} is ook
$\{(1 + 3t,\ 2 + t) : t \in \R\}$: $t = 0$ geeft $A$, $t = 1$ geeft $B$. De
parameter wegwerken ($t = y - 2$, dus $x = 1 + 3(y-2)$) levert
$x - 3y + 5 = 0$ terug.
\end{example}

\section{Onderlinge ligging van twee rechten}

\begin{proposition}[Evenwijdig of snijdend]\label{prop:g11:vect:positions}
Twee rechten met \omterm{def:g11:vect:direction}{richtingsvectoren} $\vec u$ en $\vec v$ zijn evenwijdig dan
en slechts dan als $\det(\vec u, \vec v) = 0$. In vergelijkingsvorm:
$d : ax + by + c = 0$ en $d' : a'x + b'y + c' = 0$ zijn evenwijdig dan en
slechts dan als $ab' - a'b = 0$. Niet-evenwijdige rechten ontmoeten elkaar in
precies één punt.
\end{proposition}

\begin{proof}
Evenwijdigheid betekent collineaire richtingen, en dat is
\cref{thm:g11:vect:collinearity}; met $\vec u\,(-b, a)$ en
$\vec v\,(-b', a')$ is de \omterm{def:g11:vect:det}{determinant} $(-b)a' - (-b')a = ab' - a'b$. Zijn de
rechten niet evenwijdig, dan leidt het gelijktijdig oplossen van de twee
\omterm{def:g10:algebra:equation}{vergelijkingen} (met substitutie of combinatie) tot precies één oplossing: het
stelsel gedraagt zich als een $2 \times 2$-stelsel met \omterm{def:g11:vect:det}{determinant}
verschillend van nul.
\end{proof}

\begin{method}[Snijpunt van twee rechten]\label{met:g11:vect:intersection}
Los het stelsel van de twee \omterm{def:g10:algebra:equation}{vergelijkingen} op: maak in de ene \omterm{def:g10:algebra:equation}{vergelijking}
één onbekende vrij en vul ze in de andere in, of combineer de \omterm{def:g10:algebra:equation}{vergelijkingen}
om één onbekende weg te werken. Controleer het gevonden punt altijd in
\emph{beide} \omterm{def:g10:algebra:equation}{vergelijkingen}.
\end{method}

\begin{example}\label{ex:g11:vect:intersection}
Snijd $d: x - 3y + 5 = 0$ met $d': 2x + y - 4 = 0$. Uit $d'$ volgt
$y = 4 - 2x$. Invullen in $d$ geeft
\[
x - 3(4 - 2x) + 5 = 0
\iff x - 12 + 6x + 5 = 0
\iff 7x = 7
\iff x = 1,
\]
en dan $y = 2$. De rechten ontmoeten elkaar in $(1, 2)$. Controle in $d$:
$1 - 6 + 5 = 0$.
\end{example}

\section{Oefeningen}

\begin{exercise}[$\star$]\label{exo:g11:vect:1}
Bereken bij $A(2, -1)$, $B(5, 3)$ en $C(-1, 1)$ de \omterm{def:g10:coordgeom:system}{coördinaten} van
$\vect{AB}$, $\vect{AC}$, $\vect{AB} + \vect{AC}$ en
$2\vect{AB} - 3\vect{AC}$, en het \omterm{prop:g10:coordgeom:midpoint}{midden} van $[BC]$.
\end{exercise}

\begin{exercise}[$\star$]\label{exo:g11:vect:2}
Zijn de \omterm{def:g11:vect:vector}{vectoren} \omterm{def:g11:vect:collinear}{collineair}?
\[
\vec u\,(4, -6) \text{ en } \vec v\,(-2, 3); \qquad
\vec u\,(2, 5) \text{ en } \vec v\,(3, 7); \qquad
\vec u\,(0, 3) \text{ en } \vec v\,(0, -8).
\]
\end{exercise}

\begin{exercise}[$\star$]\label{exo:g11:vect:3}
Zoek een cartesische \omterm{def:g10:algebra:equation}{vergelijking} van de rechte door $A(2, 1)$ met
\omterm{def:g11:vect:direction}{richtingsvector} $\vec u\,(3, -1)$, en daarna van de rechte door $B(0, 4)$ en
$C(2, 0)$.
\end{exercise}

\begin{exercise}[$\star$]\label{exo:g11:vect:4}
Geef voor de rechte $d: 3x - 2y + 6 = 0$ een \omterm{def:g11:vect:direction}{richtingsvector} en de
snijpunten met de twee coördinaatassen, en beslis of de punten $P(2, 6)$ en
$Q(1, 4)$ tot $d$ behoren.
\end{exercise}

\begin{exercise}[$\star\star$]\label{exo:g11:vect:5}
Liggen de punten $A(1, 2)$, $B(4, 8)$, $C(-2, -4)$ op één rechte? Dezelfde
vraag voor $D(0, 1)$, $E(2, 4)$, $F(5, 9)$.
\end{exercise}

\begin{exercise}[$\star\star$]\label{exo:g11:vect:6}
Bepaal het eventuele snijpunt:
\[
d: 2x + y - 7 = 0 \ \text{ en }\ d': x - y + 1 = 0;
\qquad
e: x - 2y + 3 = 0 \ \text{ en }\ e': -2x + 4y + 1 = 0 .
\]
\end{exercise}

\begin{exercise}[$\star\star$]\label{exo:g11:vect:7}
Geef een parametervoorstelling van de rechte $d: 2x - y + 3 = 0$, en een
cartesische \omterm{def:g10:algebra:equation}{vergelijking} van de rechte
$\bigl\{(1 + 2t,\ -3 + t) : t \in \R\bigr\}$.
\end{exercise}

\begin{exercise}[$\star\star$]\label{exo:g11:vect:8}
Zoek de \omterm{def:g10:algebra:equation}{vergelijking} van de rechte door $P(1, -2)$ die evenwijdig is met
$d: 4x - 3y + 2 = 0$, en de waarde van $m$ waarvoor de rechte
$mx + 2y - 5 = 0$ evenwijdig is met $d$.
\end{exercise}

\begin{exercise}[$\star\star$]\label{exo:g11:vect:9}
Zij $A(-1, 0)$, $B(3, 2)$ en $C(1, -4)$. Zoek een cartesische \omterm{def:g10:algebra:equation}{vergelijking}
van de \emph{zwaartelijn} van de driehoek $ABC$ uit $C$ (de rechte die $C$
met het \omterm{prop:g10:coordgeom:midpoint}{midden} van $[AB]$ verbindt).
\end{exercise}

\begin{exercise}[$\star\star$]\label{exo:g11:vect:10}
Voor welke waarde of waarden van de parameter $m$ zijn de \omterm{def:g11:vect:vector}{vectoren}
$\vec u\,(m, 2)$ en $\vec v\,(3, m - 1)$ \omterm{def:g11:vect:collinear}{collineair}?
\end{exercise}

\begin{exercise}[$\star\star\star$]\label{exo:g11:vect:11}
Zij $ABCD$ een parallellogram, $I$ het \omterm{prop:g10:coordgeom:midpoint}{midden} van $[AB]$, en $J$ het punt
bepaald door $\vect{DJ} = \frac23 \vect{DI}$. Werk in het \omterm{def:g10:coordgeom:system}{assenstelsel}
waarin $A(0,0)$, $B(1,0)$, $D(0,1)$ (zodat $C(1,1)$):
\begin{enumerate}
  \item Bereken de \omterm{def:g10:coordgeom:system}{coördinaten} van $I$ en $J$.
  \item Toon aan dat $A$, $J$ en $C$ op één rechte liggen. (Een mooi feit: de
        rechte $(DI)$ snijdt de diagonaal $(AC)$ op een derde van haar
        lengte.)
\end{enumerate}
\end{exercise}

\section{Opgave: aanvaringskoersen en slim gekozen coördinaten}

\begin{problem}[{Weekendopgave --- kruisende banen zijn nog geen aanvaring:
de kleinste nadering op zee, en stellingen bewezen door het juiste
assenstelsel te kiezen}]
\label{pb:g11:vect:1}
De rechte koersen van twee schepen kruisen elkaar --- moeten de schepen dan
botsen? Alleen wanneer ze het kruispunt \emph{op hetzelfde ogenblik}
bereiken: banen zijn verzamelingen punten, bewegingen zijn parametrisch, en
die twee verwarren heeft echte schepen doen zinken. Deze opgave houdt een
kleine scheepvaartbegeleiding op parametrische rechten
(\cref{def:g11:vect:parametric}), en wendt zich daarna tot zuivere meetkunde
met de andere superkracht van \omterm{def:g10:coordgeom:system}{coördinaten}: ze \emph{slim kiezen}, zoals in
\cref{exo:g11:vect:11}, zodat klassieke stellingen zichzelf uitrekenen.

\medskip
\noindent\textbf{Deel I --- Rechten, op drie manieren.}
\begin{enumerate}
  \item Geef een parametervoorstelling van de rechte door $A(1, 2)$ met
        \omterm{def:g11:vect:direction}{richtingsvector} $\vec u(3, -1)$, en werk daarna de parameter weg om
        een cartesische \omterm{def:g10:algebra:equation}{vergelijking} te krijgen.
  \item Lees voor de rechte $2x - 5y + 3 = 0$ een \omterm{def:g11:vect:direction}{richtingsvector} af
        (\cref{thm:g11:vect:cartesian}), zoek één punt, en schrijf een
        parametervoorstelling.
  \item Bereken het snijpunt van de twee rechten uit de vragen 1 en 2
        (\cref{met:g11:vect:intersection}).
  \item Bevestig met de \omterm{def:g11:vect:det}{determinant} (\cref{thm:g11:vect:collinearity}) dat
        die twee rechten elkaar wel moesten snijden; toon daarna aan dat
        $2x - 5y + 3 = 0$ en $-4x + 10y + 1 = 0$ strikt evenwijdig zijn.
  \item Liggen $A(2, 1)$, $B(5, 3)$, $C(11, 7)$ op één rechte
        (\cref{met:g11:vect:alignment})?
\end{enumerate}

\medskip
\noindent\textbf{Deel II --- De scheepvaartbegeleiding.}
Op tijdstip $t$ (in uren) bevindt schip $P$ zich in $(2t,\ 3 + t)$ en schip
$Q$ in $(10 - t,\ 2t - 1)$ (afstanden in zeemijl).
\begin{enumerate}[resume]
  \item Zoek de cartesische \omterm{def:g10:algebra:equation}{vergelijkingen} van de twee \emph{banen}, en
        bereken waar de banen elkaar kruisen.
  \item Op welk tijdstip passeert elk schip het kruispunt? Botsen de schepen?
        Formuleer de algemene waarschuwing in één zin: kruisende banen
        tegenover samenvallende bewegingen.
  \item Bereken het kwadraat $D^2(t)$ van de afstand tussen de schepen op
        tijdstip $t$, en herleid het tot een kwadratische uitdrukking in $t$.
        Op welk tijdstip zijn de schepen het dichtst bij elkaar, en hoe dicht
        komen ze (de toptechniek van \cref{pb:g11:quad:1})?
  \item De reglementen eisen een tussenafstand van minstens $1$ mijl. Wordt
        ze geschonden? Zo ja, los $D^2(t) < 1$ op en geef de duur van de
        overtreding.
  \item Leg uit waarom $D^2(t)$ voor \emph{elke} twee schepen op rechte
        koersen met constante snelheid altijd een \omterm{def:g11:quad:quadratic}{kwadratische functie} van
        $t$ is --- zodat de ``kleinste nadering'' altijd een topberekening is.
  \item Onderschepping: een patrouilleboot vertrekt in $t = 0$ vanuit de
        oorsprong met constante snelheid $(a, b)$ en moet schip $P$ precies
        in $t = 3$ ontmoeten. Bereken de vereiste $(a, b)$ en de snelheid van
        de patrouilleboot.
\end{enumerate}

\medskip
\noindent\textbf{Deel III --- Slim gekozen \omterm{def:g10:coordgeom:system}{coördinaten}.}
Werk in het \omterm{def:g10:coordgeom:system}{assenstelsel} van \cref{exo:g11:vect:11}: het parallellogram
$ABCD$ wordt $A(0,0)$, $B(1,0)$, $C(1,1)$, $D(0,1)$; zij $I$ het \omterm{prop:g10:coordgeom:midpoint}{midden} van
$[AB]$ en $K$ het \omterm{prop:g10:coordgeom:midpoint}{midden} van $[CD]$.
\begin{enumerate}[resume]
  \item Bereken een cartesische \omterm{def:g10:algebra:equation}{vergelijking} van de rechte $(DI)$.
  \item Snijd $(DI)$ met de diagonaal $(AC)$: haal het feit van
        \cref{exo:g11:vect:11} terug --- de snede ligt precies op een derde
        van de diagonaal.
  \item En nu de rechte $(BK)$: zoek haar \omterm{def:g10:algebra:equation}{vergelijking}, snijd met $(AC)$, en
        besluit tot de volledige klassieke stelling: \emph{de hoektransversalen
        $(DI)$ en $(BK)$ verdelen de diagonaal $(AC)$ in drie gelijke
        stukken}.
  \item Waarom was het geoorloofd de stelling in dit ene bijzondere
        \omterm{def:g10:coordgeom:system}{assenstelsel} te bewijzen? (Wat behoudt de keuze van de \omterm{def:g10:coordgeom:system}{coördinaten},
        en waarover gaat de uitspraak?) Antwoord in twee of drie zinnen.
  \item Nog een verdeling in drieën ter oefening: zij $E$ het \omterm{prop:g10:coordgeom:midpoint}{midden} van
        $[BC]$. Waar snijdt de rechte $(AE)$ de \emph{andere} diagonaal
        $(DB)$?
\end{enumerate}

\medskip
\noindent\textbf{Deel IV --- Plaatsbepalingspuzzels.}
\begin{enumerate}[resume]
  \item Gaan de drie rechten $x + y = 4$, \ $2x - y = 2$, \ $x - 2y = -2$
        door één punt? (Snijd er twee en toets de derde.)
  \item Zij $d_m$ voor elke reële $m$ de rechte $mx - y + 2 - 3m = 0$. Toon
        aan dat \emph{elke} $d_m$ door één \omterm{pb:g10:functions:1}{vast punt} gaat. (Groepeer de
        termen in $m$.)
  \item In de familie $(d_m)$: welk lid is evenwijdig met $y = 2x + 7$? En
        welk gaat door de oorsprong?
  \item Slotstuk --- de drie kostuums van een rechte: cartesisch (een test op
        lidmaatschap met één invulling), herleid (\omterm{def:g10:reffunc:affine}{richtingscoëfficiënt} en
        \omterm{def:g10:functions:graph}{grafiek} in één oogopslag), parametrisch (beweging, timing,
        onderschepping); de \omterm{def:g11:vect:det}{determinant} als orakel voor evenwijdigheid; en het
        goed gekozen \omterm{def:g10:coordgeom:system}{assenstelsel} als stellingenfabriek. Telkens één zin.
\end{enumerate}
\end{problem}
